% ID: SEL-D0037-OB-E01
% ESTADO: cotejada. Segunda pasada visual y borrador compilado revisados.
% Fuente: Selectividad/matematicas_aplicadas/Ua2MATEMÁTICAS APLICADAS examen.pdf
% SHA-256: a56739d297d4dc220d6898b24cf729dc556401a6bb292b8ca7cfb58c01c31651
% Página: 2; opción B; ejercicio original 1.

\textbf{(2.5 puntos)} Un supermercado tiene almacenados 600 kg de manzanas y 400 kg de naranjas. Para incentivar su venta elabora dos tipos de bolsas: A y B.

Las bolsas de tipo A contienen 3 kg de manzanas y 1 kg de naranjas; las bolsas de tipo B incluyen 2 kg de cada uno de los productos.

El precio de venta de la bolsa A es de 4 \texteuro{} y de 3 \texteuro{} el de la bolsa de tipo B.

Suponiendo que vende todas las bolsas preparadas, \textquestiondown{}cuántas bolsas de cada tipo debe haber elaborado para maximizar los ingresos? \textquestiondown{}A cuánto asciende el ingreso máximo?
